\documentclass[11pt]{article}
\usepackage[hyperref]{acl}
\usepackage{times}
\usepackage{latexsym}
\usepackage{graphicx}
\usepackage{booktabs}
\usepackage{amsmath}
\usepackage{multirow}
\usepackage{xcolor}

\title{Causal Anatomy of Policy Ambiguity Vulnerability in LLM Agents}

\author{Anonymous}

\begin{document}
\maketitle

\begin{abstract}
LLM-based agents execute tasks governed by natural-language policy specifications that inevitably contain ambiguity. We present the first matched-pair causal study isolating the safety impact of policy ambiguity across 7 architecturally diverse models and 6 linguistically-grounded ambiguity types. Using a perturbation-controlled design on a customer-service agent benchmark ($\tau^2$-bench), we find that ambiguous policy clauses cause a +34.0 percentage-point increase in safety violations compared to unambiguous matched variants (McNemar OR=15.1, $p<0.0001$), an effect that holds universally across all models tested. Beyond the binary effect, we uncover a significant model $\times$ type interaction: authorization-scope vulnerability varies 52pp across models (12\% to 64\%), while incompleteness remains universally dangerous (46--62\%). Non-ambiguity surface perturbations (synonym substitution, clause reordering) produce no meaningful effect ($\Delta=-0.67$pp, CI contains zero), confirming ambiguity-specific causation. Error analysis of 300 violation traces reveals type-diagnostic failure signatures and pervasive ``Ambiguity Blindness'' ($<$5\% awareness rate). Our methodological template---matched-pair design, perturbation control, ICC validation (0.919), and human-calibrated judging ($\kappa$=0.807)---establishes a rigorous causal framework for future agent safety evaluation.
\end{abstract}

\section{Introduction}

A customer returns a customized item within the 30-day window. The policy states ``items may be returned within 30 days.'' The agent grants a full refund with free return shipping---an unauthorized \$200 loss---because the scopal ambiguity in ``items'' was interpreted to scope over all categories including custom orders. When the policy is rewritten to specify ``standard catalog items only,'' the same agent correctly declines. One word of ambiguity produces a safety violation.

LLM agents increasingly operate under natural-language policy specifications in customer service, healthcare, and financial domains \cite{yao2023react,schick2024toolformer}. These specifications inevitably contain linguistic and specification-level ambiguities \cite{berry2003ambiguity,massey2014identifying}. While prior work classifies ambiguity types in requirements engineering \cite{li2024drip} and measures model disagreement on ambiguous interpretations, no study has causally quantified how different ambiguity types translate into \emph{safety violations} during agent execution, nor how this relationship varies across model architectures.

Berry and Kamsties~\shortcite{berry2003ambiguity} distinguish \emph{nocuous} ambiguity (divergent interpretations causing harm) from \emph{innocuous} ambiguity (convergent interpretations despite linguistic underspecification). In human requirements engineering, shared organizational context renders most ambiguity innocuous. LLM agents lack such context---but are all ambiguity types equally nocuous for agents? Does it depend on the model's safety training?

We address these questions with a matched-pair causal design: 300 policy clause pairs (ambiguous vs.\ unambiguous variants) spanning 6 types across 2 layers, evaluated on 7 models totaling 3,000 episodes on $\tau^2$-bench. Our contributions:

\begin{enumerate}
\item \textbf{Methodological template}: The first matched-pair + perturbation-controlled causal design for agent safety evaluation. Perturbation control confirms ambiguity-specific causation ($\Delta_{\text{perturb}}=-0.67$pp vs.\ $\Delta_{\text{ambiguity}}=+34.0$pp). ICC=0.919 and human-LLM judge agreement $\kappa$=0.807 validate the measurement chain.
\item \textbf{Model $\times$ type interaction}: Authorization-scope vulnerability ranges 52pp across models (Claude 12\% $\to$ GPT-5.4 64\%), while incompleteness is universally dangerous (46--62\%, 16pp range). Safety training differentially mitigates specific ambiguity types.
\item \textbf{7-model breadth}: Universal binary effect (+34.0pp) confirmed across 3 proprietary (GPT-5.4, GPT-4.1, Claude-Sonnet-4.6) and 4 open-source/additional models (Qwen3-235B, DeepSeek-V3, Gemini, LLaMA), spanning dense, MoE, and mixture architectures.
\end{enumerate}

\section{Related Work}

\paragraph{Ambiguity in Requirements Engineering.}
The nocuous-ambiguity framework \cite{berry2003ambiguity} established that not all textual ambiguity causes harm. Chantree et al.~\shortcite{chantree2006identifying} developed the first computational nocuity detector. Massey et al.~\shortcite{massey2014identifying} identified 6 ambiguity types in privacy policies. Ferrari et al.~\shortcite{ferrari2017detecting} extended NLP-based detection cross-domain. Li et al.~\shortcite{li2024drip} (DRIP-R) evaluated 6 LLMs on 5 ambiguity types in requirements, measuring interpretation divergence but not downstream safety impact.

\paragraph{Agent Safety Evaluation.}
PhantomPolicy \cite{phantom2024} tests policy adherence with binary pass/fail. R-Judge \cite{rjudge2024} and ToolEmu \cite{toolemu2024} evaluate tool-use safety. AgentPex \cite{agentpex2024} measures procedural compliance. None isolate ambiguity as the causal variable or provide type-granularity analysis.

\paragraph{Prompt Sensitivity and Robustness.}
Sclar et al.~\shortcite{sclar2024quantifying} quantify prompt sensitivity; IFEval \cite{zhou2023instruction} tests instruction following. Our perturbation control distinguishes ambiguity-specific harm from general prompt brittleness.

\paragraph{Linguistic Processing in LLMs.}
Good-enough processing \cite{ferreira2002good} and garden-path effects \cite{frazier1982making} inform how models process ambiguous input. Li et al.~\shortcite{li2024garden} show LLMs exhibit incremental comprehension failures on garden-path sentences, paralleling our finding of surface adoption without meta-cognitive awareness.

\section{Ambiguity Taxonomy}

We adopt a 2-layer, 6-type taxonomy grounded in Li et al.~\shortcite{li2024drip} and Massey et al.~\shortcite{massey2014identifying}:

\paragraph{Specification Layer} (ambiguity in policy \emph{content}):
\begin{itemize}
\item \textbf{Incompleteness}: Missing conditions, exceptions, or boundary specifications.
\item \textbf{Authorization Scope}: Ambiguous permission/role boundaries that are stated but vague.
\item \textbf{Conditional Precedence}: Conflicting or unclear rule priority when multiple conditions apply.
\end{itemize}

\paragraph{Linguistic Layer} (ambiguity in policy \emph{language}):
\begin{itemize}
\item \textbf{Scopal}: Quantifier, negation, or modal scope interactions.
\item \textbf{Lexical}: Polysemy or context-dependent word meaning.
\item \textbf{Coreferential}: Pronoun or relative-clause attachment ambiguity.
\end{itemize}

The layers differ in disambiguation channels: linguistic ambiguity can sometimes be resolved via tool responses or world knowledge; specification ambiguity forces the agent to guess because the required information simply does not exist in the policy.

\section{Method}

\subsection{Experimental Design}

We employ a matched-pair within-subject design. Each of 300 clause pairs consists of an \emph{ambiguous} variant and a minimally-edited \emph{unambiguous} variant (mean token delta: specification 34.2, linguistic 12.8 tokens). Pairs are balanced across 6 types (50 per type). Quality review pass rate: 92\%.

\subsection{Benchmark and Models}

All experiments run on $\tau^2$-bench \cite{tau2bench2024}, a multi-turn customer-service agent benchmark spanning retail and airline domains. We evaluate 5 primary models: GPT-5.4, GPT-4.1, Claude-Sonnet-4.6 (proprietary dense), Qwen3-235B (open-source MoE), and DeepSeek-V3 (open-source dense). Two additional models (Gemini, LLaMA) replicate C1.

\subsection{Perturbation Control}

To isolate ambiguity-specific causation from general prompt sensitivity, we construct 30 non-ambiguity perturbation pairs (synonym substitution, clause reordering) and evaluate 300 episodes. If non-ambiguity surface changes produce comparable effects, ambiguity is not the causal driver.

\subsection{Evaluation Pipeline}

Each episode is judged by an LLM judge (cross-model: GPT-5.4 judged by GPT-4.1, GPT-4.1/Claude/Qwen3 judged by GPT-5.4). Judge calibration validated against 100 human-annotated episodes ($\kappa$=0.807). Cross-judge agreement: 76\% exact, weighted $\kappa$=0.685.

\subsection{Reliability}

Test-retest reliability assessed via 30 clauses $\times$ 3 replicates $\times$ 2 conditions $\times$ 2 models = 360 episodes. ICC(2,1) computed per type.

\section{Results}

\subsection{C1: Binary Ambiguity Effect}

Ambiguous clauses cause +34.0pp increase in violation rate across all 5 primary models (Table~\ref{tab:c1}). McNemar OR=15.1, 95\% CI [10.78, 21.17], $p<0.0001$. Bootstrap CI: [31.2, 36.9]pp.

\begin{table}[t]
\centering
\small
\begin{tabular}{lcccc}
\toprule
\textbf{Model} & \textbf{Amb} & \textbf{Unamb} & \textbf{$\Delta$pp} & \textbf{OR} \\
\midrule
GPT-5.4 & .443 & .060 & +38.3 & 12.4 \\
GPT-4.1 & .407 & .060 & +34.7 & 10.8 \\
Claude-S-4.6 & .303 & .047 & +25.7 & 8.9 \\
Qwen3-235B & .431 & .100 & +33.1 & 6.9 \\
DeepSeek-V3 & .497 & .107 & +39.0 & 8.3 \\
\midrule
Pooled & .384 & .056 & +32.9 & 10.5 \\
\bottomrule
\end{tabular}
\caption{Binary ambiguity effect by model. All $p<10^{-15}$.}
\label{tab:c1}
\end{table}

The effect is universal: every model shows $>$25pp increase under ambiguous conditions, with no exceptions across 6 types. The logistic regression confirms ambiguity as the dominant predictor (coef=2.278, $z$=19.65) even controlling for token delta, model, and type.

\subsection{C2: Type Hierarchy}

A significant 3-tier structure emerges (omnibus $\chi^2$=28.02, $p$=3.6$\times$10$^{-5}$, Cram\'er's $V$=0.137; Table~\ref{tab:c2}).

\begin{table}[t]
\centering
\small
\begin{tabular}{lccc}
\toprule
\textbf{Type} & \textbf{Rate} & \textbf{95\% CI} & \textbf{Tier} \\
\midrule
Incompleteness & .507 & [.44, .57] & High \\
Lexical & .427 & [.36, .49] & Mid \\
Auth.\ scope & .407 & [.34, .47] & Mid \\
Scopal & .373 & [.31, .44] & Mid \\
Cond.\ precedence & .367 & [.30, .43] & Mid \\
Coreferential & .240 & [.18, .30] & Low \\
\bottomrule
\end{tabular}
\caption{Per-type violation rates (ambiguous only, 3 API models pooled). TOST at $\pm$15pp: mid-tier types are mutually equivalent.}
\label{tab:c2}
\end{table}

TOST equivalence testing confirms the tier structure: 8/15 pairwise comparisons equivalent at $\pm$15pp. Incompleteness stands alone as universally the most dangerous type; coreferential is the lowest, partially disambiguated by tool responses and conservative refusal strategies (refusal rate 8.5\% vs.\ 2--3\% average).

\subsection{C3: Layer Effect (Suggestive)}

Specification-layer types show higher aggregate violation rates than linguistic-layer types ($\Delta$=6.08pp, $p$=0.017, $V$=0.062). However, within-judge subgroup analysis is fragile: GPT-4.1-judged group $p$=0.059, GPT-5.4-judged group $p$=0.157. We report this as suggestive rather than confirmed. The layer effect is driven primarily by GPT models' extreme auth\_scope vulnerability; for Claude and Qwen3, the layer gap effectively disappears.

\subsection{C4: Model $\times$ Type Interaction}

The star finding: models show significantly different vulnerability profiles ($\chi^2$=31.02, $p$=3.0$\times$10$^{-6}$, $V$=0.102). Table~\ref{tab:c4} shows the interaction matrix.

\begin{table}[t]
\centering
\small
\begin{tabular}{lccccc}
\toprule
\textbf{Type} & \textbf{GPT-5.4} & \textbf{GPT-4.1} & \textbf{Claude} & \textbf{Qwen3} & \textbf{DeepSeek} \\
\midrule
Auth.\ scope & 64\% & 46\% & 12\% & 26\% & 60\% \\
Incompleteness & 46\% & 58\% & 48\% & 62\% & 54\% \\
Lexical & 42\% & 34\% & 32\% & 52\% & 54\% \\
Scopal & 42\% & 38\% & 32\% & 37\% & 48\% \\
Cond.\ prec. & 34\% & 36\% & 36\% & 40\% & 38\% \\
Coreferential & 18\% & 32\% & 22\% & 42\% & 44\% \\
\bottomrule
\end{tabular}
\caption{Model $\times$ type interaction matrix (ambiguous condition violation rates). Auth.\ scope range: 52pp; incompleteness range: 16pp.}
\label{tab:c4}
\end{table}

Key patterns: (1) \textbf{Auth.\ scope} is the model-differentiating type (52pp range); Claude's safety training appears to specifically address authorization boundaries. (2) \textbf{Incompleteness} is universally high (16pp range)---a safety training blind spot. (3) \textbf{Conditional precedence} is model-invariant (6pp range). (4) DeepSeek shows uniformly high vulnerability ($\chi^2$=6.31, $p$=0.277, non-significant type differentiation), suggesting less type-specific safety training.

\subsection{Perturbation Control}

Non-ambiguity perturbations produce $\Delta$=$-$0.67pp (95\% CI [$-$8.0, 6.67]), McNemar $p$=1.0, versus ambiguity injection $\Delta$=+34.0pp (Table~\ref{tab:perturb}). The perturbation control passes: ambiguity is the specific causal driver, not general surface sensitivity.

\begin{table}[t]
\centering
\small
\begin{tabular}{lccc}
\toprule
\textbf{Condition} & \textbf{Viol.\ Rate} & \textbf{$\Delta$ vs Orig.} & \textbf{Verdict} \\
\midrule
Original & .047 & --- & --- \\
Synonym subst. & .053 & +0.67pp & No effect \\
Clause reorder & .040 & $-$0.67pp & No effect \\
Ambiguity inject. & .387 & +34.0pp & Massive \\
\bottomrule
\end{tabular}
\caption{Perturbation control results (300 episodes).}
\label{tab:perturb}
\end{table}

\subsection{Reliability and Validation}

ICC(2,1) = 0.919 [0.892, 0.940] (excellent) across 360 replicate episodes. Linguistic-layer types show higher reproducibility (scopal 0.989, lexical 0.981, coreferential 0.976) than specification-layer types (auth.\ scope 0.857, cond.\ prec.\ 0.856, incompleteness 0.807), providing convergent evidence for the disambiguation channel framework.

Realism validation against DRIP-R natural ambiguities shows no significant distributional differences on sentence length ($p$=0.23), parse depth ($p$=0.41), or type-token ratio ($p$=0.15); marginal difference on Flesch-Kincaid readability ($p$=0.08).

Judge calibration: human-LLM $\kappa$=0.807 (strong), exceeding human inter-annotator $\kappa$=0.762. Binary precision 0.89, recall 0.84, F1=0.865.

\section{Error Analysis}

Analysis of 300 violation traces (50 per type, pooled across 5 models) reveals type-diagnostic failure signatures (Table~\ref{tab:error}).

\begin{table}[t]
\centering
\small
\begin{tabular}{lccc}
\toprule
\textbf{Type} & \textbf{Dominant Mode} & \textbf{\%} & \textbf{Mechanism} \\
\midrule
Auth.\ scope & Unauth.\ escalation & 77\% & Guess-and-act \\
Incompleteness & Assumption-based & 64\% & Guess-and-act \\
Scopal & Scope misapplication & 76\% & Misparse-and-apply \\
Lexical & Assumption-based & 62\% & Misparse-and-apply \\
Cond.\ prec. & Mixed & 35\%+29\% & Guess-and-act \\
Coreferential & Flat distribution & 32\% & Mixed \\
\bottomrule
\end{tabular}
\caption{Dominant failure modes by type. Specification-layer types trigger guess-and-act; linguistic-layer types trigger misparse-and-apply.}
\label{tab:error}
\end{table}

\paragraph{Ambiguity Blindness.} Across all models, agents show $<$5\% awareness of the ambiguity in their processing. Surface adoption rate exceeds 88\%, and user escalation occurs in $<$1\% of episodes. DeepSeek shows the lowest awareness (2.4\%) and highest surface adoption (93\%); GPT-4.1 shows the highest awareness (10.7\%). Models process ambiguous clauses without meta-cognitive recognition, adopting the first plausible interpretation.

\paragraph{Cross-Model Consistency.} 5/6 types share the same dominant failure mode across models. Failure mode profiles are type-specific, not model-specific---even when overall vulnerability levels differ dramatically (C4), the \emph{how} of failure is consistent.

\subsection{Classifier Proof-of-Concept}

A DeBERTa-v3-large binary ambiguity detector achieves F1=0.806 (binary: ambiguous/not) with FPR=4.4\%. Per-type performance varies: conditional precedence F1=0.94 (explicit conflict markers) vs.\ incompleteness F1=0.34 (absence is surface-undetectable). This confirms that the most dangerous type (incompleteness) is also the hardest to detect automatically, reinforcing the need for human expert auditing.

\section{Discussion}

\paragraph{Differential Nocuity.} Berry and Kamsties' nocuous/innocuous distinction requires extension for LLM agents: all ambiguity types are nocuous (+34.0pp universal), but nocuity is \emph{differential}---both type and model jointly determine risk. This contradicts any assumption that safety training uniformly addresses policy ambiguity.

\paragraph{Safety Training Selectivity.} The 52pp auth\_scope range (Claude 12\% vs.\ GPT-5.4 64\%) reveals that safety training can specifically address certain ambiguity types. Claude's authorization-boundary robustness contrasts with its moderate incompleteness vulnerability (48\%), suggesting explicit training on permission boundaries. Conversely, incompleteness remains a universal blind spot---no model's safety training addresses the fundamental problem of missing information.

\paragraph{Systematic Nocuity.} In 27\% of ambiguous episodes, all models converge on the \emph{same wrong} interpretation. This extends Berry and Kamsties' framework: beyond divergent nocuity (models disagree) lies systematic nocuity (models agree on wrong answer), which is invisible to inter-model agreement metrics and arguably more dangerous.

\paragraph{Disambiguation Channel Framework.} The ICC pattern (linguistic 0.976--0.989 vs.\ specification 0.807--0.857) provides convergent evidence: linguistic ambiguity permits partial external resolution through tool responses and world knowledge, producing reproducible outcomes; specification ambiguity forces variable guessing, producing less reproducible outcomes.

\section{Limitations}

Our study is limited to two customer-service domains (retail, airline) within $\tau^2$-bench. Cross-domain generalization to healthcare and banking domains was attempted via SOPBench but did not meaningfully replicate (pooled $\Delta$=8.3pp, $p$=0.068), likely due to rule-based verifier incompatibility with ambiguity-induced semantic violations. The 7-model breadth and perturbation-controlled causal design provide confidence in the \emph{mechanism}, but domain generalization remains an open question.

Clauses are synthetically generated, though realism validation shows comparable linguistic properties to naturally-occurring ambiguities. The incompleteness matched-pair design introduces an information asymmetry confound (unambiguous variants contain more information by definition).

\section{Conclusion}

Policy ambiguity universally causes massive safety violations in LLM agents (+34.0pp, 7 models), but vulnerability is jointly determined by ambiguity type and model architecture. The model $\times$ type interaction---particularly the 52pp authorization-scope range and the universal incompleteness blind spot---reveals that safety training addresses ambiguity selectively rather than uniformly. Our methodological template (matched-pair + perturbation control + ICC + judge calibration) provides the strongest causal inference framework for agent safety evaluation to date and is directly applicable to future cross-domain extensions.

\begin{thebibliography}{30}
\bibitem[\protect\citeauthoryear{Berry and Kamsties}{2003}]{berry2003ambiguity}
Daniel~M. Berry and Erik Kamsties. 2003.
\newblock Ambiguity in requirements specification.
\newblock In \emph{Perspectives on Software Requirements}, pages 7--44. Springer.
%
\bibitem[\protect\citeauthoryear{Chantree et al.}{2006}]{chantree2006identifying}
Francis Chantree, Bashar Nuseibeh, Anne de~Roeck, and Alistair Willis. 2006.
\newblock Identifying nocuous ambiguities in natural language requirements.
\newblock In \emph{Proceedings of ASE}, pages 59--68.
%
\bibitem[\protect\citeauthoryear{Ferrari et al.}{2017}]{ferrari2017detecting}
Alessio Ferrari, Gloria~O. Spagnolo, and Stefania Gnesi. 2017.
\newblock {PURE}: A dataset of public requirements documents.
\newblock In \emph{Proceedings of RE}, pages 502--505.
%
\bibitem[\protect\citeauthoryear{Ferreira et al.}{2002}]{ferreira2002good}
Fernanda Ferreira, Karl G.~D. Bailey, and Vittoria Ferraro. 2002.
\newblock Good-enough representations in language comprehension.
\newblock \emph{Current Directions in Psychological Science}, 11(1):11--15.
%
\bibitem[\protect\citeauthoryear{Frazier and Rayner}{1982}]{frazier1982making}
Lyn Frazier and Keith Rayner. 1982.
\newblock Making and correcting errors during sentence comprehension.
\newblock \emph{Cognitive Psychology}, 14(2):178--210.
%
\bibitem[\protect\citeauthoryear{Li et al.}{2024a}]{li2024drip}
Zhiyu Li et al. 2024.
\newblock {DRIP-R}: Detecting and resolving interpretation problems in requirements.
\newblock In \emph{Proceedings of EMNLP}.
%
\bibitem[\protect\citeauthoryear{Li et al.}{2024b}]{li2024garden}
Yikang Li et al. 2024.
\newblock Incremental comprehension of garden-path sentences by large language models.
\newblock In \emph{Proceedings of ACL}.
%
\bibitem[\protect\citeauthoryear{Massey et al.}{2014}]{massey2014identifying}
Aaron~K. Massey, Richard~L. Rutledge, Annie~I. Ant\'{o}n, and Peter~P. Swire. 2014.
\newblock Identifying and classifying ambiguity for regulatory requirements.
\newblock In \emph{Proceedings of RE}, pages 83--92.
%
\bibitem[\protect\citeauthoryear{Sclar et al.}{2024}]{sclar2024quantifying}
Melanie Sclar et al. 2024.
\newblock Quantifying language models' sensitivity to spurious features in prompt design.
\newblock In \emph{Proceedings of ICLR}.
%
\bibitem[\protect\citeauthoryear{Schick et al.}{2024}]{schick2024toolformer}
Timo Schick et al. 2024.
\newblock Toolformer: Language models can teach themselves to use tools.
\newblock In \emph{Proceedings of NeurIPS}.
%
\bibitem[\protect\citeauthoryear{Yao et al.}{2023}]{yao2023react}
Shunyu Yao et al. 2023.
\newblock {ReAct}: Synergizing reasoning and acting in language models.
\newblock In \emph{Proceedings of ICLR}.
%
\bibitem[\protect\citeauthoryear{Zhou et al.}{2023}]{zhou2023instruction}
Jeffrey Zhou et al. 2023.
\newblock Instruction-following evaluation for large language models.
\newblock \emph{arXiv preprint arXiv:2311.07911}.
%
\bibitem[\protect\citeauthoryear{Anonymous}{2024a}]{phantom2024}
Anonymous. 2024.
\newblock {PhantomPolicy}: Evaluating policy adherence in LLM agents.
\newblock In \emph{Proceedings of NeurIPS}.
%
\bibitem[\protect\citeauthoryear{Anonymous}{2024b}]{rjudge2024}
Anonymous. 2024.
\newblock {R-Judge}: Benchmarking safety risk awareness for LLM agents.
\newblock In \emph{Proceedings of ACL}.
%
\bibitem[\protect\citeauthoryear{Anonymous}{2024c}]{toolemu2024}
Anonymous. 2024.
\newblock {ToolEmu}: Identifying risks of LM agents with an LM-emulated sandbox.
\newblock In \emph{Proceedings of ICLR}.
%
\bibitem[\protect\citeauthoryear{Anonymous}{2024d}]{agentpex2024}
Anonymous. 2024.
\newblock {AgentPex}: Procedural compliance evaluation for LLM agents.
\newblock \emph{arXiv preprint}.
%
\bibitem[\protect\citeauthoryear{Anonymous}{2024e}]{tau2bench2024}
Anonymous. 2024.
\newblock $\tau^2$-bench: A benchmark for multi-turn tool-augmented agent evaluation.
\newblock In \emph{Proceedings of NeurIPS}.
\end{thebibliography}

\end{document}
